\documentclass[12pt]{article}
\usepackage[brazil]{babel}
\usepackage[utf8]{inputenc}
\usepackage[T1]{fontenc}
\usepackage{amsmath, amssymb}
\usepackage{geometry}
\usepackage{graphicx}
\usepackage{xcolor}
\usepackage{booktabs}
\usepackage{array}
\usepackage{fancyhdr}
\usepackage[most]{tcolorbox}

\definecolor{maincolor}{RGB}{0,100,50}
\definecolor{forestgreen}{RGB}{0,100,50}

\geometry{margin=2.5cm, headheight=53pt, top=2.5cm}
\pagestyle{fancy}
\fancyhf{}

\renewcommand{\headrulewidth}{2pt}
\renewcommand{\footrulewidth}{2pt}

\renewcommand{\headrule}{\hbox to\headwidth{%
  \color{forestgreen}\leaders\hrule height \headrulewidth\hfill}}

\renewcommand{\footrule}{\hbox to\headwidth{%
  \color{forestgreen}\leaders\hrule height \footrulewidth\hfill}}

\fancyhead[L]{%
  \includegraphics[height=1.2cm]{Universidade_de_Rio_Verde.jpg}
}

\fancyhead[R]{%
  \raggedleft
  \textcolor{forestgreen}{%
    \footnotesize
    \textbf{Propriedades da Álgebra Booleana}\\
    Me.~Sergio Souza Novak\\
    \textit{Lógica para Engenharia de Software}
  }
}

\fancyfoot[C]{%
  \textcolor{forestgreen}{\thepage}
}

\begin{document}

\begin{center}
  {\Large \textbf{\textcolor{maincolor}{Propriedades da Álgebra Booleana}}}
\end{center}

\vspace{0.5cm}

\begin{center}
\renewcommand{\arraystretch}{1.6}
\begin{tabular}{
  >{\centering\arraybackslash}m{0.5cm}
  >{\centering\arraybackslash}m{4.2cm}
  >{\centering\arraybackslash}m{4.2cm}
  >{\centering\arraybackslash}m{4.5cm}
}
\toprule
\textbf{N°} & \textbf{Forma AND (·)} & \textbf{Forma OR (+)} & \textbf{Nome da Propriedade} \\
\midrule

1 & $A \cdot 0 = 0$ & $A + 1 = 1$ & Elemento Nulo \\
\midrule

2 & $A \cdot 1 = A$ & $A + 0 = A$ & Elemento Neutro (Identidade) \\
\midrule

3 & $A \cdot A = A$ & $A + A = A$ & Idempotência \\
\midrule

4 & $A \cdot \overline{A} = 0$ & $A + \overline{A} = 1$ & Complemento \\
\midrule

5 & $A \cdot B = B \cdot A$ & $A + B = B + A$ & Comutatividade \\
\midrule

6 & $A \cdot (B \cdot C) = (A \cdot B) \cdot C$ & $A + (B + C) = (A + B) + C$ & Associatividade \\
\midrule

7 & $\overline{A \cdot B \cdot C} = \overline{A} + \overline{B} + \overline{C}$ & $\overline{A + B + C} = \overline{A} \cdot \overline{B} \cdot \overline{C}$ & Teorema de De Morgan \\
\midrule

8 & $A \cdot (B + C) = A \cdot B + A \cdot C$ & $(A+B)\cdot(A+C) = A + B \cdot C$ & Distributividade \\
\midrule

9 & $A + A \cdot B = A$ & $A \cdot (A + B) = A$ & Absorção \\
\midrule

10 & $A + \overline{A} \cdot B = A + B$ & $A \cdot (\overline{A} + B) = A \cdot B$ & Cobertura (Absorção Complementar) \\
\midrule

11 & $\overline{\overline{A}} = A$ & $\overline{\overline{A}} = A$ & Dupla Negação (Involução) \\
\midrule

12 & $A \cdot B + \overline{A} \cdot C + B \cdot C = A \cdot B + \overline{A} \cdot C$ & $(A+B) \cdot (\overline{A}+C) \cdot (B+C) = (A+B) \cdot (\overline{A}+C)$ & Teorema do Consenso \\
\midrule

13 & $C \cdot A + C \cdot \overline{A} \cdot B = C \cdot A + C \cdot B$ & $(C+A)(C+\overline{A}+B) = (C+A)(C+B)$ & Simplificação por Complemento \\

\bottomrule
\end{tabular}
\end{center}

\vspace{1cm}

\noindent\textbf{Observações:}
\begin{itemize}
  \item O símbolo $\cdot$ representa a operação AND (conjunção lógica).
  \item O símbolo $+$ representa a operação OR (disjunção lógica).
  \item O símbolo $\overline{X}$ representa a negação (NOT) da variável $X$.
  \item As propriedades valem para qualquer valor binário ($0$ ou $1$) das variáveis.
\end{itemize}

\end{document}
